% SPDX-FileCopyrightText: 2026 AntheaLiles
% SPDX-License-Identifier: CC-BY-SA-4.0
% tikz.tex — module d'export « tikz » : préambule TikZ de quickViz
% -----------------------------------------------------------------------------
% Module d'export (voir lisp/my-export-config.el) : chargé APRÈS le préambule
% de la classe (article-ua, book-ua) par le mot-clé Org
%     #+LATEX_MODULES: tikz
%
% Compilateur : LuaLaTeX. Compatible avec \DocumentMetadata{pdfstandard=ua-2,
% testphase=phase-III} : les figures produites ici portent un texte alternatif.
%
% Ce préambule ne définit aucune couleur « décorative ». Les palettes qu'il
% installe suivent les règles que le manuscrit établit :
%   — la palette qualitative fait varier la teinte le long de l'axe jaune-bleu,
%     que les déficiences les plus répandues laissent intact ;
%   — la palette séquentielle fait varier la luminance, la saturation restant
%     constante plutôt que croissante ;
%   — la palette divergente porte deux branches symétriques autour d'un neutre ;
%   — aucune information ne doit reposer sur la seule couleur : les styles de
%     trame ci-dessous servent au codage redondant.
% -----------------------------------------------------------------------------

\RequirePackage{tikz}
\RequirePackage{pgfplots}
\pgfplotsset{compat=1.18}

\usetikzlibrary{%
  arrows.meta,          % pointes de flèche paramétrables
  positioning,          % right=of, below=of
  calc,                 % coordonnées calculées
  fit,                  % nœud englobant
  backgrounds,          % plans d'arrière-plan
  shapes.geometric,     % losange des nœuds de décision
  shapes.misc,          % rectangle à coins arrondis
  patterns,             % trames nommees classiques (crosshatch, grid)
  patterns.meta,        % trames parametrables (Lines, Dots)
  decorations.pathreplacing, % accolades
  decorations.markings,
  matrix,               % grille des matrices de graphiques
  trees,                % arbre de sélection
  graphs,
  quotes}

% --- Longueur de référence -----------------------------------------------------
% Le manuscrit appelle ses figures avec width=0.9\largeurimpression. On reprend
% la même longueur si elle existe, la largeur de texte sinon.
\providecommand{\largeurimpression}{\textwidth}
\newlength{\qvlargeur}
\AtBeginDocument{\setlength{\qvlargeur}{\largeurimpression}}

% --- Palettes ------------------------------------------------------------------
% Qualitative : huit teintes séparées sur l'axe jaune-bleu, éprouvées sous
% protanopie et deutéranopie. Source : Okabe et Ito, /Color universal design/.
\definecolor{qvcat1}{HTML}{0072B2}   % bleu
\definecolor{qvcat2}{HTML}{E69F00}   % orange
\definecolor{qvcat3}{HTML}{009E73}   % vert bleuté
\definecolor{qvcat4}{HTML}{CC79A7}   % rose
\definecolor{qvcat5}{HTML}{56B4E9}   % bleu ciel
\definecolor{qvcat6}{HTML}{D55E00}   % vermillon
\definecolor{qvcat7}{HTML}{F0E442}   % jaune
\definecolor{qvcat8}{HTML}{000000}   % noir

% Séquentielle : luminance décroissante, teinte et saturation constantes.
\definecolor{qvseq1}{HTML}{DEEBF7}
\definecolor{qvseq2}{HTML}{C6DBEF}
\definecolor{qvseq3}{HTML}{9ECAE1}
\definecolor{qvseq4}{HTML}{6BAED6}
\definecolor{qvseq5}{HTML}{4292C6}
\definecolor{qvseq6}{HTML}{2171B5}
\definecolor{qvseq7}{HTML}{084594}

% Divergente : deux branches symétriques autour d'un neutre clair.
\definecolor{qvdivA3}{HTML}{8C510A}
\definecolor{qvdivA2}{HTML}{BF812D}
\definecolor{qvdivA1}{HTML}{DFC27D}
\definecolor{qvdiv0}{HTML}{F5F5F5}
\definecolor{qvdivB1}{HTML}{80CDC1}
\definecolor{qvdivB2}{HTML}{35978F}
\definecolor{qvdivB3}{HTML}{01665E}

% Mise en évidence : un accent, le reste en gris. La palette d'évidence du
% manuscrit ne compte que deux valeurs — ce qu'on regarde, et le reste.
\definecolor{qvaccent}{HTML}{D55E00}
\definecolor{qvfond}{HTML}{BDBDBD}

% Encre et repères : jamais du noir pur pour les traits de structure.
\definecolor{qvencre}{HTML}{1A1A1A}
\definecolor{qvgrille}{HTML}{D9D9D9}
\definecolor{qvaxe}{HTML}{595959}
\definecolor{qvalerte}{HTML}{B2182B}

% --- Codage redondant ----------------------------------------------------------
% Une trame double la couleur partout où l'information ne doit pas dépendre
% d'elle seule. Elle demande des aplats assez larges pour que sa densité se voie.
% Une trame ne remplace pas la couleur, elle s'y ajoute : TikZ ne pose qu'un
% remplissage par chemin, d'où le style « fond », qui peint l'aplat avant la
% trame. L'usage est donc : [qv barre, fond=qvcat1, trame A=white].
\tikzset{
  fond/.style={preaction={fill=#1}},
  fond/.default=white,
  trame A/.style={pattern={Lines[angle=45,distance=3pt,line width=.6pt]},
                  pattern color=#1},
  trame A/.default=qvencre,
  trame B/.style={pattern={Lines[angle=-45,distance=3pt,line width=.6pt]},
                  pattern color=#1},
  trame B/.default=qvencre,
  trame C/.style={pattern={Dots[distance=3pt,radius=.7pt]}, pattern color=#1},
  trame C/.default=qvencre,
  trame D/.style={pattern=crosshatch, pattern color=#1},
  trame D/.default=qvencre,
  trame vide/.style={fill=white},
}

% --- Texte alternatif ----------------------------------------------------------
% Une figure TikZ n'est pas décrite par elle-même. L'environnement qvfigure
% prend le texte alternatif en premier argument obligatoire et le pose sur la
% structure de la figure lorsque tagpdf est actif ; sinon il l'ignore
% silencieusement, de sorte que le document compile aussi hors mode balisé.
\newcommand{\qvalt}[1]{%
  \ifdefined\tagpdfsetup
    \tagpdfsetup{tabsorder=structure}%
    \tagstructbegin{tag=Figure,alt={#1}}\tagmcbegin{artifact}%
  \fi}
\newcommand{\qvaltfin}{%
  \ifdefined\tagpdfsetup
    \tagmcend\tagstructend
  \fi}

% #1 (optionnel) : options passées à tikzpicture. #2 : texte alternatif.
\NewDocumentEnvironment{qvfigure}{O{} m}
  {\qvalt{#2}\begin{tikzpicture}[qv, #1]}
  {\end{tikzpicture}\qvaltfin}

% --- Style de base -------------------------------------------------------------
% Tout dessin de ce manuscrit s'ouvre sur le style « qv ». Il fixe la police, la
% graisse des traits et la couleur d'encre, de sorte qu'aucune figure n'ait à
% les redéclarer.
\tikzset{
  qv/.style={
    line width=.7pt,
    draw=qvencre,
    text=qvencre,
    font=\footnotesize\sffamily,
    >={Stealth[length=5pt,width=4pt]},
    every node/.append style={inner sep=3pt},
  },
}

% --- Arbre de sélection --------------------------------------------------------
% Les nœuds de l'arbre de la figure « Parcours de sélection ». La question est un
% losange, la condition un rectangle, la feuille — un type de graphique — un
% rectangle à coins arrondis. La forme double la couleur : l'arbre reste lisible
% en noir et blanc.
\tikzset{
  qv question/.style={
    diamond, aspect=2.2, draw=qvaxe, fill=qvseq1,
    align=center, inner sep=2pt, minimum width=26mm},
  qv condition/.style={
    rectangle, draw=qvaxe, fill=white,
    align=center, minimum width=24mm, minimum height=7mm},
  qv feuille/.style={
    rounded corners=2pt, draw=qvencre, fill=qvseq2,
    align=center, minimum width=24mm, minimum height=7mm},
  qv feuille repli/.style={
    qv feuille, fill=white, dashed},
  qv branche/.style={draw=qvaxe, -{Stealth[length=4pt]}, rounded corners=3pt},
  qv etiquette/.style={
    font=\scriptsize\sffamily\itshape, fill=white, inner sep=1.5pt,
    midway, sloped=false},
}

% --- Schémas de graphiques -----------------------------------------------------
% Les fiches illustrent des principes de tracé, non des données réelles. Ces
% styles servent à dessiner un graphique schématique : deux axes, une grille
% discrète, des marques.
\tikzset{
  qv axe/.style={draw=qvaxe, line width=.6pt, -{Stealth[length=4pt,width=3pt]}},
  qv grille/.style={draw=qvgrille, line width=.4pt},
  qv reperes/.style={draw=qvaxe, line width=.5pt},
  qv barre/.style={draw=qvencre, line width=.5pt, fill=qvcat1},
  qv barre fond/.style={qv barre, fill=qvfond},
  qv barre accent/.style={qv barre, fill=qvaccent},
  qv ligne/.style={draw=qvcat1, line width=1.1pt, line cap=round,
                   line join=round},
  qv ligne second/.style={qv ligne, draw=qvcat2, dash pattern=on 3pt off 2pt},
  qv ligne tiers/.style={qv ligne, draw=qvcat3, dash pattern=on 1pt off 1.6pt},
  qv point/.style={circle, fill=qvcat1, draw=none, inner sep=0pt,
                   minimum size=3pt},
  qv aire/.style={fill=qvseq3, draw=qvcat1, line width=.7pt, fill opacity=.55},
  qv zero/.style={draw=qvencre, line width=.9pt},
  qv annotation/.style={font=\scriptsize\sffamily, text=qvaxe, align=left},
  qv etiquette directe/.style={font=\scriptsize\sffamily, inner sep=1.5pt,
                               anchor=west},
  qv alerte/.style={draw=qvalerte, line width=.8pt, dash pattern=on 2pt off 2pt},
  qv rupture/.style={decorate,
    decoration={zigzag, amplitude=.8pt, segment length=3pt}},
}

% Un cadre d'axes complet en une commande : \qvcadre{largeur}{hauteur}
\newcommand{\qvcadre}[2]{%
  \draw[qv axe] (0,0) -- (#1,0);
  \draw[qv axe] (0,0) -- (0,#2);}

% --- Matrice de graphiques -----------------------------------------------------
% Les panneaux d'une matrice partagent leurs bornes et leurs axes. Ce style ne
% pose qu'un titre d'ordonnées par ligne et un titre d'abscisses par colonne.
\tikzset{
  qv panneau/.style={draw=qvgrille, fill=white, line width=.5pt,
                     minimum width=22mm, minimum height=16mm, anchor=south west},
  qv titre panneau/.style={font=\scriptsize\sffamily, anchor=north west,
                           text=qvaxe},
  qv matrice/.style={matrix of nodes, nodes in empty cells,
                     column sep=2mm, row sep=2mm, nodes={qv panneau}},
}

% --- Zone de prescription et repères de lecture --------------------------------
\tikzset{
  qv zone/.style={fill=qvseq1, draw=none, rounded corners=2pt},
  qv zone alerte/.style={fill=qvalerte!8, draw=qvalerte, dash pattern=on 2pt off 2pt,
                         rounded corners=2pt},
  qv accolade/.style={decorate, decoration={brace, amplitude=4pt, raise=2pt},
                      draw=qvaxe, line width=.5pt},
  qv legende/.style={draw=qvgrille, fill=white, rounded corners=2pt,
                     inner sep=4pt, font=\scriptsize\sffamily},
  qv puce/.style={rectangle, minimum width=7pt, minimum height=7pt,
                  inner sep=0pt, draw=qvencre, line width=.4pt},
}

% --- Chemin de décision --------------------------------------------------------
% Ce que le moteur expose : la forme retenue, les formes écartées, et le motif de
% l'écart. La forme écartée se barre et porte sa contrainte violée ; elle n'est
% pas seulement grisée, faute de quoi le motif ne serait pas lisible.
\tikzset{
  qv retenue/.style={qv feuille, fill=qvseq2, draw=qvencre, line width=1pt},
  qv ecartee/.style={qv feuille, fill=white, draw=qvfond, text=qvaxe},
  qv motif/.style={font=\scriptsize\sffamily\itshape, text=qvalerte,
                   anchor=west, align=left},
  qv egalite/.style={draw=qvaxe, line width=.5pt, dash pattern=on 1pt off 1.5pt},
}

% --- Pgfplots : accord avec les règles du manuscrit ----------------------------
% L'échelle des barres part de zéro, la grille reste discrète, et les séries
% portent une marque distincte en plus de leur couleur.
\pgfplotsset{
  qv/.style={
    width=.9\qvlargeur, height=.55\qvlargeur,
    axis lines=left,
    axis line style={draw=qvaxe, line width=.6pt},
    tick align=outside, tick style={draw=qvaxe, line width=.5pt},
    every tick label/.append style={font=\scriptsize\sffamily, text=qvaxe},
    label style={font=\footnotesize\sffamily},
    title style={font=\footnotesize\sffamily\bfseries, align=left},
    grid=major, grid style={draw=qvgrille, line width=.4pt},
    % La légende se pose hors du cadre pour ne jamais recouvrir une marque.
    % L'étiquetage direct reste préférable sur une figure imprimée : voir le
    % composant « Annotation » du manuscrit.
    legend style={at={(1.02,1)}, anchor=north west,
                  draw=qvgrille, fill=white, font=\scriptsize\sffamily,
                  cells={anchor=west}},
    legend cell align=left,
    cycle list={%
      {qvcat1, mark=*, mark size=1.6pt},
      {qvcat2, mark=square*, mark size=1.5pt, dash pattern=on 3pt off 2pt},
      {qvcat3, mark=triangle*, mark size=1.9pt, dash pattern=on 1pt off 1.6pt},
      {qvcat4, mark=diamond*, mark size=1.8pt, dash pattern=on 4pt off 1.5pt},
      {qvcat6, mark=pentagon*, mark size=1.7pt, dash pattern=on 2pt off 2pt on 5pt off 2pt},
    },
  },
  qv barres/.style={qv, ybar, ymin=0, bar width=8pt,
                    every axis plot/.append style={fill=qvcat1, draw=qvencre,
                                                   line width=.4pt}},
}
